% ============================================================
%  Documento di Progetto - Informatica con Laboratorio 6CFU
%  Equazione del Calore 2D e Riordinamento nelle
%  Fattorizzazioni Cholesky
%
%  Compilare con:  pdflatex design.tex  (due passate)
% ============================================================

\documentclass[a4paper,12pt]{article}

% ---- Pacchetti base ----------------------------------------
\usepackage[T1]{fontenc}
\usepackage[utf8]{inputenc}
\usepackage[italian]{babel}
\usepackage{lmodern}

% ---- Layout ------------------------------------------------
\usepackage[
  top=2.0cm, bottom=2.0cm,
  left=2.0cm, right=2.0cm,
  headheight=15pt
]{geometry}
\usepackage{parskip}
\usepackage{microtype}

% ---- Matematica --------------------------------------------
\usepackage{amsmath, amssymb, amsthm}

% ---- Tabelle e box -----------------------------------------
\usepackage{booktabs}
\usepackage{array}

% ---- Pseudocodice ------------------------------------------
\usepackage{listings}
\lstset{
  basicstyle=\footnotesize\ttfamily,
  breaklines=true,
  breakatwhitespace=true,
  numbers=left,
  numberstyle=\tiny\color{gray},
  stepnumber=1,
  numbersep=10pt,
  frame=lines,
  aboveskip=10pt,
  belowskip=10pt,
  xleftmargin=15pt,
  escapeinside={(*}{*)}
}

% ---- Float e TikZ ------------------------------------------
\usepackage{float}
\usepackage{graphicx}
\usepackage{tikz}
\usetikzlibrary{shapes.geometric, arrows.meta, positioning, fit, backgrounds}

% ---- Colori ------------------------------------------------
\usepackage{xcolor}

% ---- Link --------------------------------------------------
\usepackage[hidelinks]{hyperref}

% ---- Intestazioni ------------------------------------------
\usepackage{fancyhdr}
\pagestyle{fancy}
\fancyhf{}
\rhead{\small Equazione del Calore 2D -- Documento di Progetto}
\lhead{\small Giosuè Aiello}
\cfoot{\thepage}

% ---- Titolo ------------------------------------------------
\title{%
  \LARGE\textbf{Equazione del Calore 2D}\\[6pt]
  \LARGE\textbf{e Riordinamento nelle Fattorizzazioni Cholesky}\\[12pt]
  \large Documento di Progetto\\[4pt]
  \normalsize Informatica con Laboratorio -- A.A. 2025/2026
}
\author{Giosuè Aiello}
\date{}

% ============================================================
\begin{document}
\maketitle
\thispagestyle{empty}
% ============================================================

% ------------------------------------------------------------
\section{Impostazione del problema}
% ------------------------------------------------------------

L'obiettivo del progetto è la risoluzione numerica dell'equazione del calore in regime stazionario in due dimensioni spaziali. Il dominio di calcolo è il quadrato unitario $\Omega = [0,1]^2$, e si cerca la funzione di temperatura $u : \Omega \to \mathbb{R}$ che soddisfa l'equazione ellittica:

\begin{equation}
  \kappa \left(\frac{\partial^2 u}{\partial x^2} + \frac{\partial^2 u}{\partial y^2}\right) + f(x,y) = 0,
  \qquad (x,y) \in \Omega,
  \label{eq:calore}
\end{equation}

con condizioni al contorno di Dirichlet sul bordo $\partial\Omega$:
\[
  u(x,y) = u_0(x,y), \qquad (x,y) \in \partial\Omega.
\]

I parametri fisici del problema sono: coefficiente di diffusività termica $\kappa = 0.01$ e termine sorgente $f(x,y) = e^{-10(x^2+y^2)}$, che modella una sorgente di calore concentrata in prossimità dell'origine con decadimento gaussiano. Nel caso base si considera $u_0 \equiv 0$ (condizioni omogenee).

\subsection*{Discretizzazione alle differenze finite e appiattimento}

Il dominio continuo $\Omega$ viene sostituito da una griglia cartesiana uniforme composta da $(N+2) \times (N+2)$ nodi fisici, spaziati tra loro di un passo costante $h = 1/(N+1)$. I singoli nodi sono identificati dalle coordinate spaziali discrete $(x_i, y_j) = (i \cdot h, j \cdot h)$, dove gli indici spaziali $i, j$ variano nell'intervallo $\{0, 1, \dots, N+1\}$. 

In questo reticolo si distinguono due categorie di nodi:
\begin{itemize}
    \item \textbf{Nodi di bordo:} corrispondono agli indici $i \in \{0, N+1\}$ oppure $j \in \{0, N+1\}$. In questi punti la temperatura è nota a priori grazie alle condizioni di Dirichlet e non costituisce un'incognita del sistema.
    \item \textbf{Nodi interni:} corrispondono agli indici $1 \le i \le N$ e $1 \le j \le N$. Questi rappresentano i veri gradi di libertà del problema. Indichiamo con $u_{i,j} \approx u(x_i, y_j)$ l'approssimazione numerica della temperatura in tali punti. Essi costituiscono le $N^2$ incognite del nostro sistema.
\end{itemize}

Poiché i risolutori algebrici di sistemi lineari operano su array monodimensionali, è strettamente necessario ``appiattire'' (flattening) la griglia 2D in un vettore unidimensionale $\mathbf{u} \in \mathbb{R}^{N^2}$. A tale scopo, introduciamo una mappa di indicizzazione biunivoca $\phi(i,j)$, definita come l'ordinamento naturale riga per riga:
\begin{equation}
    k = \phi(i,j) = (i-1) \cdot N + (j-1), \qquad k \in \{0, 1, \dots, N^2-1\}.
\end{equation}
Tramite questa funzione, la temperatura bidimensionale $u_{i,j}$ viene mappata in modo univoco nell'elemento $k$-esimo del vettore delle incognite $\mathbf{u}$. Analogamente, il termine noto e i contributi noti al bordo andranno a comporre il vettore globale $\mathbf{b} \in \mathbb{R}^{N^2}$.

\begin{figure}[htbp]
    \centering
    \begin{tikzpicture}[scale=0.9]
        \colorlet{bdraw}{black!60}
        \colorlet{bfill}{gray!20}
        \colorlet{idraw}{blue!80!black}
        \colorlet{ifill}{blue!70!white}
        \colorlet{rdraw}{red!80!black}
        \colorlet{rfill}{red!70!white}

        \fill[blue!4] (1,1) rectangle (6,6);
        \draw[black!35, thin] (0,0) grid (7,7);

        \tikzset{
            bnode/.style={circle, draw=bdraw, fill=bfill, inner sep=1.5pt, minimum size=5pt},
            inode/.style={circle, draw=idraw, fill=ifill, inner sep=1.5pt, minimum size=5pt},
            rnode/.style={circle, draw=rdraw, fill=rfill, inner sep=2pt, minimum size=7pt}
        }

        \foreach \x in {0,...,7} { \node[bnode] at (\x,0) {}; \node[bnode] at (\x,7) {}; }
        \foreach \y in {1,...,6} { \node[bnode] at (0,\y) {}; \node[bnode] at (7,\y) {}; }
        \foreach \x in {1,...,6} {
            \foreach \y in {1,...,6} { \node[inode] (N-\x-\y) at (\x,\y) {}; }
        }

        \node[rnode] (R) at (3,3) {};
        \draw[->, >=stealth, rdraw, line width=1.2pt] (R) -- (N-2-3) node[midway, below, font=\tiny\bfseries, rdraw] {$u_{i-1,j}$};
        \draw[->, >=stealth, rdraw, line width=1.2pt] (R) -- (N-4-3) node[midway, below, font=\tiny\bfseries, rdraw] {$u_{i+1,j}$};
        \draw[->, >=stealth, rdraw, line width=1.2pt] (R) -- (N-3-2) node[midway, right, font=\tiny\bfseries, rdraw] {$u_{i,j-1}$};
        \draw[->, >=stealth, rdraw, line width=1.2pt] (R) -- (N-3-4) node[midway, right, font=\tiny\bfseries, rdraw] {$u_{i,j+1}$};

        \node[above right=3pt, text=rdraw, fill=white, inner sep=1.5pt, font=\small\bfseries] at (R) {$(x_i, y_j)$};
        \node[below left=2pt, font=\footnotesize] at (0,0) {$(0,0)$};
        \node[below right=2pt, font=\footnotesize] at (7,0) {$(1,0)$};
        \node[above right=2pt, font=\footnotesize] at (7,7) {$(1,1)$};
        \draw[->, thick, black!60] (-0.3,0) -- (7.5,0) node[right,font=\small]{$x$};
        \draw[->, thick, black!60] (0,-0.3) -- (0,7.5) node[above,font=\small]{$y$};
        \draw[<->, gray!70, thin] (0,-0.55) -- (1,-0.55) node[midway, below, font=\tiny]{$h$};
    \end{tikzpicture}
    \caption{Griglia di discretizzazione: nodi di bordo (grigio), nodi interni (blu) e i quattro vicini dello stencil per il nodo $(x_i, y_j)$ (frecce rosse).}
\end{figure}

Approssimando le derivate seconde con lo stencil a cinque punti, l'equazione~\eqref{eq:calore} valutata nel punto interno $(x_i, y_j)$ diventa:

\begin{equation}
  \frac{\kappa}{h^2}\bigl(u_{i-1,j} + u_{i+1,j} + u_{i,j-1} + u_{i,j+1} - 4\,u_{i,j}\bigr) + f(x_i,y_j) = 0.
  \label{eq:stencil}
\end{equation}

L'insieme delle equazioni~\eqref{eq:stencil}, combinate tramite l'appiattimento $\phi(i,j)$, costituisce il sistema lineare algebrico globale $A\mathbf{u} = \mathbf{b}$. In questa formulazione matriciale:
\begin{itemize}
    \item $\mathbf{u} \in \mathbb{R}^{N^2}$ è il \textbf{vettore delle incognite}, la cui componente $k$-esima contiene la temperatura approssimata $u_{i,j}$ nel nodo interno $(x_i, y_j)$ mappato dall'indice $k$.
    \item $\mathbf{b} \in \mathbb{R}^{N^2}$ è il \textbf{vettore dei termini noti}, la cui componente $k$-esima raccoglie il termine sorgente (cambiato di segno, $-f(x_i, y_j)$) unito agli eventuali contributi noti derivanti dai nodi di bordo adiacenti al nodo $(i,j)$.
    \item $A \in \mathbb{R}^{N^2 \times N^2}$ è la \textbf{matrice dei coefficienti}, che risulta sparsa (avendo al più 5 elementi non nulli per riga, in corrispondenza del pattern dello stencil), simmetrica e definita negativa.
\end{itemize}

\subsection*{Fattorizzazione di Cholesky e fill-in}

Poiché $-A$ è simmetrica definita positiva (SPD), si calcola la fattorizzazione di Cholesky $-A = LL^T$ con $L$ triangolare inferiore sparsa. Durante la fattorizzazione si generano elementi non nulli in posizioni che erano zero in $-A$: questo fenomeno è detto \emph{fill-in} e dipende fortemente dall'ordine con cui le incognite sono enumerate.

Con l'ordinamento naturale per righe la larghezza di banda di $A$ è $\Theta(N)$ e il fill-in di $L$ (che determina il \textbf{costo spaziale} o occupazione di memoria) cresce come $O(N^3)$. La tecnica di \emph{nested dissection} riduce questo costo spaziale a $O(N^2 \log N)$, con un risparmio sostanziale di memoria per valori grandi di $N$.

\subsection*{Nested dissection: approccio geometrico}

La struttura di $A$ è codificata dal \emph{grafo di adiacenza} $G = (V, E)$: ogni nodo $v \in V$ corrisponde a una variabile $u_{i,j}$ e ogni arco $(u,v) \in E$ indica che $A_{uv} \neq 0$. Per la griglia regolare:
\[
  (x_i, y_j) \leftrightarrow (x_k, y_l) \iff (i=k,\, |j-l|=1) \;\text{ oppure }\; (j=l,\, |i-k||=1).
\]

La \emph{nested dissection} individua ricorsivamente un separatore $V_S$ che divide $V$ in due insiemi bilanciati $V_1$ e $V_2$ senza archi diretti tra loro. L'ordinamento risultante elenca prima $V_1$, poi $V_2$, infine $V_S$, in modo che il separatore appaia in fondo alla matrice, limitando la propagazione del fill-in.

In questo progetto si adotta un \textbf{approccio geometrico puro}: il separatore non è determinato esplorando il grafo di adiacenza, bensì è definito direttamente tramite le coordinate spaziali dei nodi. Scelto l'asse di taglio $c \in \{x, y\}$ e calcolato l'indice mediano $\hat{c}$, si pone:
\[
  V_1 = \{u_{i,j}\mid c < \hat{c}\},\quad
  V_2 = \{u_{i,j}\mid c > \hat{c}\},\quad
  V_S = \{u_{i,j}\mid c = \hat{c}\}.
\]
Il procedimento viene applicato ricorsivamente a $V_1$ e $V_2$, alternando il taglio lungo $x$ e $y$. Per griglie cartesiane regolari, il piano mediano è \emph{al contempo} un separatore geometrico e un separatore topologico del grafo: non esistono archi di adiacenza tra $V_1$ e $V_2$ (i punti sono separati da almeno un passo reticolare lungo l'asse di taglio), quindi l'approccio geometrico produce esattamente il medesimo separatore ottimale che si otterrebbe esplorando la topologia del grafo, senza il costo di tale esplorazione.

\subsubsection*{Approccio geometrico vs.\ approccio su grafo}

Esistono due famiglie di algoritmi per la nested dissection:

\begin{enumerate}
  \item \textbf{Approccio su grafo} (es.\ bisection spettrale, METIS): esplora esplicitamente il grafo di adiacenza $G=(V,E)$ e calcola autovettori del Laplaciano. Richiede $O(|E|) = O(N^2)$ solo per costruire e visitare il grafo, più il costo delle euristiche di partizionamento (generalmente $O(N^2 \alpha(N))$). Tali metodi sono necessari per griglie non strutturate o domini irregolari.
  
  \item \textbf{Approccio geometrico} (adottato in questo progetto): ignora completamente la lista di adiacenza. Il separatore viene individuato tramite il \emph{piano mediano} sull'asse corrente: dati gli indici interi $(i, j)$ dei nodi, si calcola $\hat{c} = \lfloor(\min_c + \max_c)/2\rfloor$ e si partiziona $V_S = \{v \mid c_v = \hat{c}\}$, $V_1 = \{v \mid c_v < \hat{c}\}$, $V_2 = \{v \mid c_v > \hat{c}\}$.
\end{enumerate}

\paragraph{Perché l'approccio geometrico è superiore per griglie regolari.}
Per una griglia cartesiana uniforme $N \times N$ l'approccio geometrico è \emph{equivalente} a quello su grafo in termini di qualità del separatore, ma \emph{superiore} sotto ogni aspetto implementativo e in parte anche dei costi:

\begin{itemize}
  \item \textbf{Correttezza del separatore}: il piano mediano $c = \hat{c}$ divide la griglia in due metà con al massimo $\lceil N/2 \rceil$ nodi ciascuna. Non esistono archi di adiacenza tra $V_1$ e $V_2$ (i nodi adiacenti si trovano a distanza reticolare $h$ lungo l'asse di taglio, quindi il nodo con $c = \hat{c}$ funge da buffer). Il separatore geometrico \emph{coincide} esattamente con il separatore topologico del grafo.
  \item \textbf{Complessità}: $\mathcal{O}(N^2 \log N)$ deterministico, senza euristiche. L'approccio su grafo richiede almeno $O(N^2)$ per la sola lettura degli archi, più il costo dell'algoritmo di partizionamento.
  \item \textbf{Efficienza in cache}: accesso sequenziale al solo vettore di nodi \\ (\texttt{coords.txt} $\rightarrow$ \texttt{std::vector<Node>}), nessuna lista di puntatori per adiacenze.
  \item \textbf{Robustezza}: utilizzo di confronti su indici interi ($(i,j)$) anziché coordinate floating-point, eliminando i rischi di arrotondamento.
\end{itemize}

% ------------------------------------------------------------
\newpage
\section{Architettura del progetto}
% ------------------------------------------------------------

\subsection*{Modulo principale: \texttt{HES--HeatEquationSolver}}

Il sistema viene concepito come un unico risolutore logico denominato \texttt{HES--HeatEquationSolver}, articolato in quattro sotto-moduli specializzati. I tre moduli in C++ comunicano tra loro attraverso file di testo con formato fisso che fungono da interfaccia verificabile indipendentemente; il modulo Python costituisce il punto di ingresso per la fase numerica finale.

\medskip
\begin{center}
\begin{tikzpicture}[
  node distance = 0.65cm and 2.0cm,
  modcpp/.style = {
    rectangle, rounded corners=3pt,
    draw=blue!60!black, fill=blue!8, line width=0.8pt,
    text width=4.6cm, align=center,
    font=\small\ttfamily\bfseries, minimum height=1.1cm, inner sep=5pt
  },
  modpy/.style = {
    rectangle, rounded corners=3pt,
    draw=orange!70!black, fill=orange!10, line width=0.8pt,
    text width=4.6cm, align=center,
    font=\small\ttfamily\bfseries, minimum height=1.1cm, inner sep=5pt
  },
  filebox/.style = {
    rectangle,
    draw=black!40, fill=gray!5, line width=0.5pt,
    text width=3.4cm, align=center,
    font=\footnotesize\ttfamily, minimum height=0.80cm, inner sep=3pt
  },
  startnode/.style = {
    rectangle, rounded corners=3pt,
    draw=black!70, fill=white, line width=0.8pt,
    font=\small\bfseries, align=center,
    minimum height=0.85cm, minimum width=2.5cm
  },
  arr/.style     = {-{Stealth[length=6pt]}, line width=0.9pt, black!70},
  arrfile/.style = {-{Stealth[length=5pt]}, line width=0.6pt,
                    dashed, dash pattern=on 4pt off 2pt, black!50}
]

\node[startnode] (IN) {Input $N$};
\node[modcpp] (T1) [below=0.65cm of IN]
  {task1\_grid.cpp\\[2pt]{\normalfont\footnotesize\rmfamily Griglia e grafo}};
\node[modcpp] (T2) [below=1.2cm of T1]
  {task2\_reorder.cpp\\[2pt]{\normalfont\footnotesize\rmfamily Nested Dissection}};
\node[modcpp] (T3) [below=1.2cm of T2]
  {task3\_assemble.cpp\\[2pt]{\normalfont\footnotesize\rmfamily Sistema lineare}};
\node[modpy]  (NB) [below=1.2cm of T3]
  {solver.ipynb\\[2pt]{\normalfont\footnotesize\rmfamily Cholesky $+$ analisi}};

\node[filebox] (F1) [right=2.0cm of T1] {coords.txt\\connectivity.txt};
\node[filebox] (F2) [right=2.0cm of T2] {ordering.txt};
\node[filebox] (F3) [right=2.0cm of T3] {A.txt \quad rhs.txt};

\draw[arr] (IN.south) -- (T1.north);
\draw[arr] (T1.south) -- (T2.north);
\draw[arr] (T2.south) -- (T3.north);
\draw[arr] (T3.south) -- (NB.north);

\draw[arrfile] (T1.east) -- (F1.west);
\draw[arrfile] (T2.east) -- (F2.west);
\draw[arrfile] (T3.east) -- (F3.west);

% Task2 legge SOLO coords.txt (approccio geometrico puro)
\draw[arrfile] (F1.south) -- ++(0,-0.20) -| ([xshift=1.0cm]T2.east)
  node[midway, above, font=\tiny\itshape, text=black!50] {coords.txt} -- (T2.east);
\draw[arrfile] (F2.south) -- ++(0,-0.35) -| ([xshift=1.0cm]T3.east) -- (T3.east);
\draw[arrfile] (F3.south) -- ++(0,-0.35) -| ([xshift=1.0cm]NB.east)  -- (NB.east);

\end{tikzpicture}

\smallskip
{\footnotesize\rmfamily
\rule{1.5em}{0.8pt}\quad Flusso di esecuzione\qquad
\raisebox{2pt}{\tikz{\draw[dashed, dash pattern=on 4pt off 2pt] (0,0) -- (1.5em,0);}}\quad File di interfaccia\qquad
{\color{blue!60!black}$\blacksquare$}\quad Moduli C++\qquad
{\color{orange!70!black}$\blacksquare$}\quad Modulo Python
}
\end{center}

\bigskip

\noindent Le strutture dati fondamentali dell'intero sistema sono:

\begin{itemize}
  \item \textbf{Liste di adiacenza} (\texttt{std::vector<std::vector<int>>}) nel Modulo 1 per la costruzione e scrittura del grafo di adiacenza.
  \item \textbf{Vettore di nodi logici} (\texttt{std::vector<Node>}) nel Modulo 2: \texttt{struct Node} contiene \texttt{id, i, j}; le coordinate in virgola mobile vengono lette e ignorate dal parser per minimizzare l'occupazione in cache. Questo consente la nested dissection geometrica calcolando il mediano sugli indici senza caricare in memoria la lista di adiacenza.
  \item \textbf{Ricorsione classica} nel Modulo 2: l'algoritmo divide-et-impera è implementato ricorsivamente come da specifiche, potendo contare su una profondità limitata dell'albero delle chiamate.
  \item \textbf{Vettori di permutazione} (\texttt{std::vector<int>} \texttt{perm}, \texttt{inv\_perm}) per la mappatura bidirezionale tra ordinamento naturale e riordinato.
  \item \textbf{Formato sparso CSC} (\texttt{scipy.sparse.csc\_matrix}) per la matrice $A$ nel modulo Python, richiesto nativamente da CHOLMOD.
\end{itemize}

% ------------------------------------------------------------
\newpage
\section{Specifiche dei moduli}
% ------------------------------------------------------------

I quattro sotto-moduli vengono descritti con livello di dettaglio crescente. Per ciascuno si indicano input, output, funzionalità richieste, strutture dati, invarianti, dipendenze e complessità attesa.

% ---- Modulo 1 -----------------------------------------------
\subsection*{Modulo 1 -- \texttt{task1\_grid.cpp}: generazione della griglia e del grafo}

Il Modulo 1 è responsabile della creazione della struttura geometrica del dominio e del relativo grafo di adiacenza. Costituisce il punto di ingresso dell'intera pipeline e non dipende da nessun altro modulo.

\subsubsection*{Formalizzazione matematica}

Il dominio $\Omega = [0,1]^2$ viene discretizzato con passo $h = \frac{1}{N+1}$. L'insieme dei nodi interni (le incognite) è:
\[
  V = \{ (x_i, y_j) \mid x_i = i \cdot h,\; y_j = j \cdot h,\quad i, j \in \{1, \dots, N\} \}, \qquad |V| = N^2.
\]
Il grafo di adiacenza $G = (V, E)$ ha archi:
\[
  E = \left\{ \{u, v\} \subset V \;\middle|\; \sqrt{(x_i - x_k)^2 + (y_j - y_l)^2} = h \right\}.
\]
Ogni nodo interno ha al massimo 4 vicini (sinistra, destra, basso, alto). A ogni nodo viene assegnato un identificatore intero univoco tramite la mappa $\phi(i,j) = (i-1) \cdot N + (j-1)$.

\subsubsection*{Specifiche}

\begin{itemize}
  \item \textbf{Input:} Intero $N > 0$ passato da riga di comando; se assente o non valido, si attiva un prompt iterativo con opzione di uscita (\texttt{q}).
  \item \textbf{Output:}
    \begin{itemize}
      \item \texttt{output/coords.txt} -- una riga per nodo: \texttt{id i j x y}
      \item \texttt{output/connectivity.txt} -- una riga per arco: \texttt{e u v} (dove \texttt{e} è l'indice progressivo dell'arco e $u < v$)
    \end{itemize}
  \item \textbf{Funzionalità:}
    \begin{itemize}
      \item Parsing e validazione di $N$ con fallback interattivo a ciclo (loop).
      \item Generazione dei $N^2$ nodi con ID progressivo e coordinate spaziali (precisione \texttt{double}).
      \item Costruzione della lista di adiacenza bidirezionale e scrittura degli archi senza duplicati ($u < v$).
    \end{itemize}
  \item \textbf{Strutture dati:} \texttt{std::vector<std::vector<int>>} di dimensione $N^2$, pre-allocata con \texttt{reserve(4)} per ciascun nodo.
  \item \textbf{Invarianti:}
    \begin{itemize}
      \item La lista di adiacenza è sempre simmetrica: se $v \in \text{adj}[u]$ allora $u \in \textit{adj}[v]$.
      \item Ogni arco appare esattamente una volta in \texttt{connectivity.txt} (condizione $u < v$).
      \item Il numero di righe in \texttt{coords.txt} è esattamente $N^2$.
    \end{itemize}
  \item \textbf{Dipendenze:} Nessuna. Questo modulo è il primo della pipeline.
  \item \textbf{Complessità attesa:} $\mathcal{O}(N^2)$ in tempo e spazio.
\end{itemize}

\subsubsection*{Logica implementativa (pseudo-codice)}
\begin{minipage}{\linewidth}
\begin{lstlisting}
Programma Principale (Generazione Griglia):
  -- Inizializzazione --
  Leggi la Dimensione_Griglia (N) passata dall'utente
  Se la dimensione non e' valida:
    Richiedi l'inserimento ciclicamente finche' non e' un numero positivo

  -- Creazione dei Nodi e degli Archi --
  Per ogni Riga (i) da 1 fino a N:
    Per ogni Colonna (j) da 1 fino a N:
      Calcola l'Identificativo_Univoco del nodo attuale
      Calcola le coordinate spaziali reali (X, Y) del nodo
      
      Salva i dati del nodo (ID, riga, colonna, X, Y) nel file 'coords.txt'
      
      -- Connessioni (Grafo) --
      Se il nodo non e' sull'ultimo bordo destro:
        Crea un arco di collegamento con il nodo successivo a destra
      Se il nodo non e' sull'ultimo bordo inferiore:
        Crea un arco di collegamento con il nodo sottostante

  -- Salvataggio della Connettivita' --
  Per ogni nodo e per ogni suo collegamento:
    Se non e' un duplicato (evita andata e ritorno):
      Salva il collegamento nel file 'connectivity.txt'
\end{lstlisting}
\end{minipage}

\clearpage
% ---- Modulo 2 -----------------------------------------------
\subsection*{Modulo 2 -- \texttt{task2\_reorder.cpp}: nested dissection geometrica}

Il Modulo 2 costituisce il cuore algoritmico del progetto. Calcola una permutazione ottimizzata dei nodi tramite la \textit{nested dissection geometrica}, al fine di minimizzare il fill-in durante la fattorizzazione di Cholesky.

\subsubsection*{Formalizzazione matematica}

L'algoritmo riceve in input il vettore $S$ di nodi con coordinate intere $(i,j)$ e reali $(x,y)$, e individua ricorsivamente:
\begin{align*}
  V_1 &= \{ u \in S \mid c_u < \hat{c} \} \\
  V_2 &= \{ u \in S \mid c_u > \hat{c} \} \\
  V_S &= \{ u \in S \mid c_u = \hat{c} \}
\end{align*}
dove $c \in \{i, j\}$ è l'indice dell'asse di taglio e $\hat{c} = \lfloor (\min_{u \in S} c_u + \max_{u \in S} c_u)/2 \rfloor$. L'ordinamento risultante è $V_1 \to V_2 \to V_S$, ricorsivo con alternanza degli assi.

\subsubsection*{Specifiche}

\begin{itemize}
  \item \textbf{Input:} \texttt{output/coords.txt} (prodotto dal Modulo~1). Il file \texttt{connectivity.txt} \emph{non viene letto}: per l'approccio geometrico la struttura topologica del grafo è ridondante.
  \item \textbf{Output:} \texttt{output/ordering.txt} -- una riga per posizione: \texttt{new\_id old\_id}
  \item \textbf{Funzionalità:}
    \begin{itemize}
      \item Lettura delle sole coordinate da \texttt{coords.txt}.
      \item Partizionamento tramite ricorsione classica (funzione che richiama se stessa).
      \item Calcolo del piano mediano con indici interi, con alternanza degli assi $i$/$j$ ad ogni livello.
      \item Verifica dell'invariante: \texttt{ordering.size() == N\textsuperscript{2}}.
      \item Salvataggio dell'ordinamento su \texttt{ordering.txt}.
    \end{itemize}
  \item \textbf{Strutture dati:} \texttt{struct Node \{id, i, j\}}; \texttt{std::vector<int> ordering}.
  \item \textbf{Invarianti:}
    \begin{itemize}
      \item Il vettore \texttt{ordering} è una permutazione valida di $\{0, \dots, N^2-1\}$.
      \item I nodi di $V_S$ compaiono sempre \emph{dopo} i nodi di $V_1$ e $V_2$ alla stessa profondità.
      \item L'asse di taglio si alterna ad ogni livello (0 = $i$, 1 = $j$).
    \end{itemize}
  \item \textbf{Dipendenze:} Modulo~1 (solo \texttt{coords.txt}).
  \item \textbf{Complessità attesa:} $\mathcal{O}(N^2 \log N)$ in tempo; $\mathcal{O}(N^2)$ in spazio.
\end{itemize}

\clearpage
\begin{figure}[p]
\centering

% --- PRIMO ELEMENTO (IN ALTO) ---
\begin{tikzpicture}[
  scale=0.95, transform shape,
  node distance = 0.5cm and 0.8cm,
  startend/.style = {rectangle, rounded corners=6pt, draw=black!70, fill=gray!15, minimum width=2cm, minimum height=0.6cm, text centered, font=\sffamily\footnotesize\bfseries},
  process/.style = {rectangle, draw=blue!60!black, fill=blue!8, minimum width=3cm, minimum height=0.6cm, text centered, text width=5.5cm, font=\sffamily\footnotesize},
  decision/.style = {diamond, aspect=2.0, draw=red!70!black, fill=red!10, text centered, font=\sffamily\footnotesize, inner sep=1pt, minimum width=2.5cm},
  arrow/.style = {thick, -{Stealth[length=6pt]}},
  looparrow/.style = {thick, dashed, -{Stealth[length=6pt]}, draw=purple!80!black}
]

% Nodi
\node (start) [startend, text width=4cm] {Avvia partizionamento:\\\texttt{Dissezione(Nodi, Direzione)}};
\node (dec1) [decision, below=of start] {Zero nodi?};
\node (ret1) [startend, right=0.8cm of dec1] {Termina (Ritorna)};

\node (dec2) [decision, below=of dec1] {Un solo nodo?};
\node (add1) [process, right=0.8cm of dec2, text width=3.5cm] {Salva il nodo nel\\Nuovo Ordinamento};
\node (ret2) [startend, right=0.5cm of add1] {Termina};

\node (calc) [process, below=of dec2] {Calcola la coordinata Mediana lungo la Direzione corrente};
\node (part) [process, below=of calc] {Dividi i Nodi in:\\ \textbf{Prima Meta'} ($<$ Mediana)\\ \textbf{Seconda Meta'} ($>$ Mediana)\\ \textbf{Separatore} ($=$ Mediana)};
\node (rec1) [process, below=of part] {Ricorsione: Applica Dissezione sulla Prima Meta'};
\node (rec2) [process, below=of rec1] {Ricorsione: Applica Dissezione sulla Seconda Meta'};
\node (addS) [process, below=of rec2] {Accoda tutti i nodi del Separatore\\al Nuovo Ordinamento (Post-ordine)};
\node (end) [startend, below=of addS] {Fine del ramo corrente};

% Archi
\draw [arrow] (start) -- (dec1);
\draw [arrow] (dec1) -- node[above, font=\sffamily\scriptsize] {Sì} (ret1);
\draw [arrow] (dec1) -- node[left, font=\sffamily\scriptsize] {No} (dec2);
\draw [arrow] (dec2) -- node[above, font=\sffamily\scriptsize] {Sì} (add1);
\draw [arrow] (add1) -- (ret2);
\draw [arrow] (dec2) -- node[left, font=\sffamily\scriptsize] {No} (calc);
\draw [arrow] (calc) -- (part);
\draw [arrow] (part) -- (rec1);
\draw [arrow] (rec1) -- (rec2);
\draw [arrow] (rec2) -- (addS);
\draw [arrow] (addS) -- (end);

% Frecce di iterazione / Ricorsione (Linee Ortogonali)
\draw [looparrow] (rec1.west) -- ++(-1.5,0) |- node[pos=0.25, left, font=\sffamily\scriptsize, text=purple!80!black] {Nuova Iterazione} (start.west);
\draw [looparrow] (rec2.west) -- ++(-2.0,0) |- node[pos=0.25, left, font=\sffamily\scriptsize, text=purple!80!black] {Nuova Iterazione} (start.west);

\end{tikzpicture}
\caption{Diagramma di flusso (\texttt{task2}).}
\label{fig:nd_flowchart}

\vspace{1.5cm}

% --- SECONDO ELEMENTO (IN BASSO) ---
\begin{tikzpicture}[scale=0.8, every node/.style={circle, draw, minimum size=0.6cm, inner sep=0.5pt, font=\scriptsize\bfseries}]
  % 1. Bottom-Left Block (1..9)
  \foreach \x in {1,2,3} {
    \foreach \y in {1,2,3} {
      \pgfmathtruncatemacro{\idx}{(\y-1)*3 + \x}
      \node[fill=blue!10, draw=blue!50] (N-\x-\y) at (\x,\y) {\idx};
    }
  }
  % 2. Top-Left Block (10..18)
  \foreach \x in {1,2,3} {
    \foreach \y in {5,6,7} {
      \pgfmathtruncatemacro{\idx}{(\y-5)*3 + \x + 9}
      \node[fill=blue!10, draw=blue!50] (N-\x-\y) at (\x,\y) {\idx};
    }
  }
  % 3. Bottom-Right Block (19..27)
  \foreach \x in {5,6,7} {
    \foreach \y in {1,2,3} {
      \pgfmathtruncatemacro{\idx}{(\y-1)*3 + (\x-4) + 18}
      \node[fill=green!10, draw=green!50] (N-\x-\y) at (\x,\y) {\idx};
    }
  }
  % 4. Top-Right Block (28..36)
  \foreach \x in {5,6,7} {
    \foreach \y in {5,6,7} {
      \pgfmathtruncatemacro{\idx}{(\y-5)*3 + (\x-4) + 27}
      \node[fill=green!10, draw=green!50] (N-\x-\y) at (\x,\y) {\idx};
    }
  }
  % 5. Left Horizontal Separator (37..39)
  \foreach \x in {1,2,3} {
    \pgfmathtruncatemacro{\idx}{\x + 36}
    \node[fill=cyan!35, draw=cyan!80, thick] (N-\x-4) at (\x,4) {\idx};
  }
  % 6. Right Horizontal Separator (40..42)
  \foreach \x in {5,6,7} {
    \pgfmathtruncatemacro{\idx}{(\x-4) + 39}
    \node[fill=yellow!45, draw=yellow!80, thick] (N-\x-4) at (\x,4) {\idx};
  }
  % 7. Primary Vertical Separator (43..49)
  \foreach \y in {1,...,7} {
    \pgfmathtruncatemacro{\idx}{\y + 42}
    \node[fill=red!30, draw=red!80, thick] (N-4-\y) at (4,\y) {\idx};
  }
  % Disegna gli archi della griglia
  \foreach \x in {1,...,7} {
    \foreach \y in {1,...,6} {
      \pgfmathtruncatemacro{\ynext}{\y+1}
      \draw[gray!60, thin] (N-\x-\y) -- (N-\x-\ynext);
    }
  }
  \foreach \y in {1,...,7} {
    \foreach \x in {1,...,6} {
      \pgfmathtruncatemacro{\xnext}{\x+1}
      \draw[gray!60, thin] (N-\x-\y) -- (N-\xnext-\y);
    }
  }
\end{tikzpicture}
\vspace{0.4cm}
\caption{Esempio di ordinamento su griglia.}
\label{fig:nd_matrix}
\end{figure}
\clearpage

\subsubsection*{Logica implementativa (pseudo-codice)}
\begin{minipage}{\linewidth}
\begin{lstlisting}
Programma Principale:
  1. Carica l'elenco dei nodi dal file 'coords.txt'
  2. Avvia la partizione: Nested_Dissection(Tutti_i_nodi, Asse_Orizzontale)
  3. Verifica che tutti i nodi siano stati riordinati correttamente
  4. Salva la permutazione finale nel file 'ordering.txt'

Procedura Ricorsiva Nested_Dissection(Sottoinsieme_Nodi, Asse_Corrente):
  -- Casi base (condizioni di arresto) --
  Se il Sottoinsieme_Nodi e' vuoto:
    Interrompi e ritorna
  Se il Sottoinsieme_Nodi contiene un solo nodo:
    Aggiungi l'identificativo del nodo al nuovo ordinamento
    Interrompi e ritorna

  -- Individuazione del separatore --
  Trova la coordinata minima e massima lungo l'Asse_Corrente
  Calcola il Punto_Medio tra la minima e la massima

  -- Partizionamento geometrico --
  Suddividi i nodi in tre nuovi raggruppamenti:
    - Prima_Meta (nodi con coordinata minore del Punto_Medio)
    - Seconda_Meta (nodi con coordinata maggiore del Punto_Medio)
    - Separatore (nodi posizionati esattamente sul Punto_Medio)

  -- Alternanza della direzione di taglio (da x a y, e viceversa) --
  Prossimo_Asse = Alterna_Direzione(Asse_Corrente)

  -- Esplorazione ricorsiva (Divide et Impera) --
  Richiama Nested_Dissection(Prima_Meta, Prossimo_Asse)
  Richiama Nested_Dissection(Seconda_Meta, Prossimo_Asse)

  -- Post-ordine (fill-in minimizzato) --
  Per ogni nodo contenuto nel Separatore:
    Aggiungi l'identificativo del nodo al nuovo ordinamento
\end{lstlisting}
\end{minipage}

\clearpage
% ---- Modulo 3 -----------------------------------------------
\subsection*{Modulo 3 -- \texttt{task3\_assemble.cpp}: assemblaggio del sistema lineare}

Il Modulo 3 costruisce il sistema lineare $Au = b$ che approssima l'equazione differenziale. Supporta sia l'ordinamento naturale che quello ottimizzato via il flag \texttt{-r}, e implementa la parte opzionale delle condizioni al bordo non omogenee.

\subsubsection*{Formalizzazione matematica}

Applicando lo stencil a 5 punti con passo $h = 1/(N+1)$ e diffusività $\kappa$, i coefficienti della matrice per ogni nodo $k$ (nel nuovo ordinamento) sono:
\begin{align*}
  A_{k,k}  &= -\frac{4\kappa}{h^2} \\
  A_{k,k'} &= +\frac{\kappa}{h^2} \qquad \text{per ogni vicino interno } k'
\end{align*}
Il termine noto è $b_k = -f(x_{i(k)}, y_{j(k)})$. Se il vicino $k'$ cade sul bordo (condizione di Dirichlet), il suo valore $u_0$ è noto e viene incorporato nel termine noto:
\begin{equation*}
  b_k \;\mathrel{{-}{=}}\; \frac{\kappa}{h^2} \cdot u_0\!\left(x_{\text{bordo}},\, y_{\text{bordo}}\right).
\end{equation*}

\begin{figure}[htbp]
    \centering
    \begin{tikzpicture}[scale=1.3]
        \tikzset{
            cell/.style={rectangle, draw=gray!50, fill=gray!6, minimum size=1.1cm, font=\small},
            cen/.style={rectangle, draw=red!70!black, fill=red!10, line width=1.2pt, minimum size=1.1cm, font=\small\bfseries},
            nbr/.style={rectangle, draw=blue!60!black, fill=blue!8, line width=0.8pt, minimum size=1.1cm, font=\small}
        }
        \foreach \c in {-2,-1,0,1,2} {
            \foreach \r in {-2,-1,0,1,2} { \node[cell] at (\c*1.15, \r*1.15) {}; }
        }
        \node[nbr] (top) at (0, 1.15)  {$+\frac{\kappa}{h^2}$};
        \node[nbr] (bot) at (0,-1.15)  {$+\frac{\kappa}{h^2}$};
        \node[nbr] (lft) at (-1.15,0)  {$+\frac{\kappa}{h^2}$};
        \node[nbr] (rgt) at ( 1.15,0)  {$+\frac{\kappa}{h^2}$};
        \node[cen] (cen) at (0,0)      {$-\frac{4\kappa}{h^2}$};

        \node[font=\footnotesize, blue!70!black, above=2pt] at (top.north) {$u_{i,j+1}$};
        \node[font=\footnotesize, blue!70!black, below=2pt] at (bot.south) {$u_{i,j-1}$};
        \node[font=\footnotesize, blue!70!black, left=2pt]  at (lft.west)  {$u_{i-1,j}$};
        \node[font=\footnotesize, blue!70!black, right=2pt] at (rgt.east)  {$u_{i+1,j}$};
        \node[font=\footnotesize, red!70!black, below=2pt]  at (cen.south) {$u_{i,j}$};

        \draw[->, >=stealth, red!60, thick] (cen) -- (top);
        \draw[->, >=stealth, red!60, thick] (cen) -- (bot);
        \draw[->, >=stealth, red!60, thick] (cen) -- (lft);
        \draw[->, >=stealth, red!60, thick] (cen) -- (rgt);
        
        \node[font=\normalsize] at (2.0*1.15+0.9, 0) {$= -f(x_i,y_j)$};
    \end{tikzpicture}
    \caption{Maschera dello stencil a 5 punti: bilancio dell'equazione differenziale per l'assemblaggio di una riga di $A$.}
\end{figure}

\subsubsection*{Specifiche}

\begin{itemize}
  \item \textbf{Input:} Intero $N$ e flag opzionale \texttt{-r} da riga di comando; \texttt{output/coords.txt}; \texttt{output/ordering.txt} (solo se \texttt{-r}).
  \item \textbf{Output:} \texttt{output/A.txt} (triplette \texttt{i j val}) e \texttt{output/rhs.txt} (una riga per valore di $b$).
  \item \textbf{Funzionalità:}
    \begin{itemize}
      \item Lettura di \texttt{coords.txt} e costruzione della mappa $\phi^{-1}: (i,j) \to \text{id}$.
      \item Se \texttt{-r}: lettura di \texttt{ordering.txt} e costruzione di \texttt{perm} e \texttt{inv\_perm}; altrimenti vettori identità.
      \item Menu interattivo per la scelta della condizione al bordo $u_0$ (selezione avviene una sola volta, prima del ciclo di assemblaggio).
      \item Assemblaggio dello stencil con gestione delle condizioni di Dirichlet.
      \item Scrittura di \texttt{A.txt} e \texttt{rhs.txt} con precisione floating-point a 10 cifre decimali.
    \end{itemize}
  \item \textbf{Strutture dati:} \texttt{std::vector<int> perm, inv\_perm} di dimensione $N^2$; \texttt{struct Triplet \{int row, col; double val\}}; \texttt{std::vector<Triplet>} pre-allocato con \texttt{reserve(5*N*N)}; \texttt{std::vector<double> b}.
  \item \textbf{Invarianti:}
    \begin{itemize}
      \item \texttt{perm} e \texttt{inv\_perm} sono permutazioni inverse l'una dell'altra: \texttt{inv\_perm[perm[k]] == k} per ogni $k$.
      \item La matrice $A$ è simmetrica: se la tripletta $(k, k', v)$ è presente, lo è anche $(k', k, v)$.
      \item Il numero di righe in \texttt{rhs.txt} è esattamente $N^2$.
    \end{itemize}
  \item \textbf{Dipendenze:} Modulo 1 (\texttt{coords.txt}); Modulo 2 (\texttt{ordering.txt}, solo se \texttt{-r}).
  \item \textbf{Complessità attesa:} $\mathcal{O}(N^2)$ in tempo (al più 5 entrate per nodo) e spazio.
\end{itemize}

\subsubsection*{Logica implementativa (pseudo-codice)}
\begin{minipage}{\linewidth}
\begin{lstlisting}
Programma Principale (Costruzione del Sistema Lineare):
  -- Inizializzazione --
  Leggi i parametri di avvio: Dimensione Griglia, Uso del Riordinamento, Tipo Condizioni Bordo
  Carica in memoria tutte le informazioni sui nodi create dal Modulo 1

  -- Configurazione dell'Ordine --
  Se e' attivo il Riordinamento (Nested Dissection):
    Carica il Nuovo Ordinamento e costruisci le mappe di conversione (Vecchia <-> Nuova posizione)
  Altrimenti:
    Usa l'ordinamento classico riga per riga

  -- Costruzione delle Equazioni --
  Per ogni riga della matrice globale (che corrisponde alla Nuova Posizione del nodo):
    Individua chi era originariamente questo nodo e recupera le sue coordinate (X, Y)
    
    Salva il termine centrale dell'equazione (l'interazione del nodo con se stesso)
    Calcola il calore generato internamente (f(x,y)) e mettilo nel vettore dei termini noti
    
    -- Interazioni con i vicini --
    Per ogni direzione possibile (Sopra, Sotto, Sinistra, Destra):
      Calcola la posizione del nodo vicino in quella direzione
      
      Se il vicino esce dal bordo della griglia:
        Applica la temperatura del bordo (Condizione di Dirichlet) spostandola nel termine noto
      Altrimenti (il vicino e' interno alla griglia):
        Scopri qual e' la Nuova Posizione di questo vicino
        Salva il termine di interazione termica tra il nodo attuale e questo vicino
        
  -- Salvataggio --
  Salva tutte le interazioni calcolate nel file 'A.txt'
  Salva tutti i termini noti nel file 'rhs.txt'
\end{lstlisting}
\end{minipage}

\subsubsection*{Condizioni al bordo disponibili (parte opzionale)}

\begin{center}
\begin{tabular}{clp{7cm}}
\toprule
\textbf{Scelta} & \textbf{Funzione $u_0(x,y)$} & \textbf{Utilità di verifica} \\
\midrule
0 & $0$ & Caso base omogeneo (default) \\
1 & $10$ & Con $f=0$, la soluzione deve essere $\equiv 10$ \\
2 & $x$ & Con $f=0$, la soluzione interna deve coincidere con $x$ \\
3 & $\sin(\pi x)$ & Bordo continuo; si annulla agli angoli \\
4 & $x^2 - y^2$ & $\nabla^2(x^2-y^2)=0$: soluzione numerica esatta \\
5 & $\mathbf{1}_{y=1}$ & Shock termico: solo il bordo superiore è caldo \\
\bottomrule
\end{tabular}
\end{center}

\clearpage
% ---- Modulo 4 (Task 4 e Task 5) -----------------------------
\subsection*{Modulo 4 -- \texttt{solver.ipynb}: risoluzione (Task 4) e analisi (Task 5)}

Il Modulo 4 è il punto terminale della pipeline. In termini di specifiche progettuali, questo modulo accorpa coerentemente sia le richieste della \textbf{Task 4} (costruzione della matrice sparsa CSC e fattorizzazione di Cholesky) sia della \textbf{Task 5} (benchmark per $N \in \{32, \dots, 1024\}$, \textit{spy plot} della struttura e grafici della soluzione). 

Si è optato per questa scelta architetturale poiché entrambe le task si collocano naturalmente nello stesso ambiente di esecuzione (Python). Piuttosto che frammentare la logica in script separati, un unico Jupyter Notebook (\texttt{solver.ipynb}) funge da orchestratore finale: legge i dati generati dal Modulo 3, esegue i calcoli numerici e restituisce contestualmente l'analisi visiva e prestazionale.

\subsubsection*{Specifiche}

\begin{itemize}
  \item \textbf{Input:}
    \begin{itemize}
      \item \texttt{output/A.txt} e \texttt{output/rhs.txt} (prodotti dal Modulo 3).
      \item Parametro $N$ (o serie $N \in \{32, 64, 128, 256, 512, 1024\}$ per il benchmark).
    \end{itemize}
  \item \textbf{Output:}
    \begin{itemize}
      \item Vettore soluzione $u \in \mathbb{R}^{N^2}$, rimappato sulla griglia originale per la visualizzazione.
      \item Spy plot della matrice $A$ e del fattore $L$ per ordinamento naturale e nested dissection.
      \item Mappa di calore 2D e superficie 3D della soluzione $u(x,y)$.
      \item Grafico log-log del fill-in \texttt{nnz(L)} e dei tempi di fattorizzazione al variare di $N$.
    \end{itemize}
  \item \textbf{Funzionalità:}
    \begin{itemize}
      \item Lettura di \texttt{A.txt} e costruzione di \texttt{scipy.sparse.coo\_matrix}, con conversione in \texttt{csc\_matrix}.
      \item Fattorizzazione $-A = LL^T$ tramite \texttt{sksparse.cholmod.cholesky} con \texttt{order="natural"}.
      \item Risoluzione del sistema con \texttt{scipy.sparse.linalg.spsolve\_triangular} (forward e backward substitution).
      \item Automazione dell'intera pipeline C++ tramite \texttt{subprocess} per ciascun valore di $N$ e ordinamento.
      \item Misurazione dei tempi con \texttt{time.perf\_counter}.
    \end{itemize}
  \item \textbf{Strutture dati:}
    \begin{itemize}
      \item \texttt{scipy.sparse.coo\_matrix} per l'assemblaggio iniziale; \texttt{csc\_matrix} per la fattorizzazione.
      \item \texttt{numpy.ndarray} per la soluzione $u$ e il termine noto $b$.
      \item Oggetto \texttt{Factor} di CHOLMOD per estrarre la matrice $L$ esplicita.
    \end{itemize}
  \item \textbf{Invarianti:}
    \begin{itemize}
      \item La matrice costruita da \texttt{coo\_matrix} deve essere simmetrica e definita negativa.
      \item Il vettore soluzione $u$ deve soddisfare $\|Au - b\|_2 / \|b\|_2 < 10^{-10}$ (residuo relativo). Questa severa soglia assicura l'assenza di cancellazioni catastrofiche prodotte dall'accumulo di errori in virgola mobile, tollerando al contempo il fisiologico malcondizionamento $\mathcal{O}(N^2)$ della matrice Laplaciana all'aumentare di $N$.
    \end{itemize}
  \item \textbf{Dipendenze:} Modulo 3 (\texttt{A.txt}, \texttt{rhs.txt}); librerie \texttt{scipy}, \texttt{numpy}, \texttt{scikit-sparse}, \texttt{matplotlib}.
  \item \textbf{Complessità attesa:} Fill-in $O(N^3)$ con ordinamento naturale; $O(N^2 \log N)$ con nested dissection. Tempo di fattorizzazione dominato dal fill-in.
\end{itemize}

\subsubsection*{Logica implementativa (pseudo-codice)}
\begin{minipage}{\linewidth}
\begin{lstlisting}
Programma Principale (Orchestratore in Python):

  -- FASE 1: Risoluzione e Validazione Fisica (es. N=32) --
  Avvia pipeline C++ (Modulo 1, Modulo 3)
  Leggi la matrice sparsa A da 'A.txt' e il termine noto b da 'rhs.txt'
  Fattorizza A = L L^T (Cholesky)
  Risolvi il sistema lineare (Sostituzione in Discesa e in Risalita)
  Costruisci la griglia 2D includendo le condizioni al bordo
  Mostra la mappa di calore e la superficie 3D delle temperature
  
  -- FASE 2: Benchmark Prestazionale Asintotico --
  Per ogni Dimensione N da testare (32, 64, 128, ...):
    Per ogni Ordinamento (Naturale vs Nested Dissection):
      
      Avvia pipeline C++ per lo specifico N e Ordinamento
      Leggi ESCLUSIVAMENTE la matrice A da 'A.txt'
      
      Fai partire il Cronometro
      Fattorizza A = L L^T (Cholesky)
      Ferma il Cronometro
      
      Registra il Tempo di Fattorizzazione e conta il Fill-in (Numeri Diversi da Zero in L)
      
  -- Visualizzazione Finale Benchmark --
  Disegna i grafici comparativi Log-Log (Tempo e Fill-in rispetto a N^2)
\end{lstlisting}
\end{minipage}

% --\section{Analisi di Complessità Asintotica ($\mathcal{O}$)}
% ------------------------------------------------------------

La sostenibilità dell'intera pipeline di calcolo si regge sull'efficienza dei tre moduli preprocessori scritti in C++. Affinché il solutore numerico Python finale possa gestire griglie estremamente dense (come $N=1024$), è imperativo che la costruzione del sistema lineare e il calcolo del riordinamento non costituiscano essi stessi un collo di bottiglia.

\vspace{0.3cm}

\subsection{Il collo di bottiglia dell'Ordinamento Naturale}
L'approccio ingenuo (ordinamento naturale lessicografico) fallisce irrimediabilmente per esaurimento della memoria (OOM - Out of Memory) a dimensioni elevate. Poiché la larghezza di banda della matrice $A$ è pari a $N$, il fattore di Cholesky $L$ soffre di un fill-in (riempimento di zeri con valori non nulli) proporzionale a:
\begin{equation}
    \text{Fill-in (Naturale)} = \mathcal{O}(N^3)
\end{equation}
Per $N=1024$, il numero di incognite è $V \approx 10^6$. Il fattore $L$ naturale richiederebbe di immagazzinare $\approx 10^9$ elementi in virgola mobile a doppia precisione (circa 8 GB di pura RAM solo per la matrice, ignorando l'overhead), rendendo la fattorizzazione instabile o impossibile sui computer portatili standard. La \textit{Nested Dissection} riduce drasticamente questo fill-in a $\mathcal{O}(N^2 \log N)$, rendendo il problema computazionalmente trattabile.

\vspace{0.3cm}

La tabella seguente riassume la complessità asintotica in Spazio e Tempo dei nostri algoritmi C++ (con $V = N^2$ incognite totali):

\begin{table}[H]
\centering
\begin{tabular}{lll}
\toprule
\textbf{Modulo (C++)} & \textbf{Complessità Temporale} & \textbf{Complessità Spaziale} \\
\midrule
\texttt{task1\_grid} & $\mathcal{O}(N^2)$ & $\mathcal{O}(N^2)$ \\
\texttt{task2\_reorder} & $\mathcal{O}(N^2 \log N)$ & $\mathcal{O}(N^2)$ \\
\texttt{task3\_assemble} & $\mathcal{O}(N^2)$ & $\mathcal{O}(N^2)$ \\
\bottomrule
\end{tabular}
\end{table}

\subsection{Modulo 1: Generazione della Topologia in $\mathcal{O}(N^2)$}
Il problema fisico è impostato su un dominio quadrato bidimensionale discretizzato con una griglia cartesiana. Le incognite numeriche (i gradi di libertà interni) sono esattamente $V = N^2$. 

L'algoritmo di instradamento effettua una singola scansione spaziale tramite due cicli \texttt{for} annidati (iterando lungo gli assi $y$ e $x$). Ad ogni iterazione, l'elaboratore esegue un set costante di operazioni logiche con costo $\mathcal{O}(1)$:
\begin{itemize}
    \item Calcolo dell'identificativo unidimensionale \texttt{id}.
    \item Assegnazione delle coordinate reali $x$ e $y$.
    \item Mappatura dei vicini (creazione degli archi).
\end{itemize}
Poiché vengono eseguiti passaggi in tempo costante per ciascuno degli $N^2$ nodi, il tempo globale scala linearmente con il numero totale di incognite, $\mathcal{O}(N^2)$.
Lo spazio in memoria RAM è dominato dall'allocazione delle strutture dati primarie (come \texttt{std::vector<Node>}), dimensionate esattamente a $N^2$ istanze. Ciò garantisce una limitazione spaziale rigorosa e stabile a $\mathcal{O}(N^2)$.

\vspace{0.5cm}

\subsection{Modulo 2: Nested Dissection Geometrica in $\mathcal{O}(N^2 \log N)$}
Il modulo di riordinamento sfrutta la disposizione geometrica della griglia anziché esplorare il grafo. Adotta un approccio ricorsivo di tipo \textit{Divide et Impera} per dividere il dominio in sotto-problemi.

Questa suddivisione genera un albero di decisione con profondità logaritmica pari a $\mathcal{O}(\log V)$, dove $V = N^2$ è il numero totale di incognite. Poiché a ogni livello dell'albero il lavoro complessivo svolto su tutti i sotto-problemi è proporzionale a $V$, la complessità temporale complessiva risulta pari a:
\begin{equation}
    T(V) = \mathcal{O}(V \log V) \implies \mathcal{O}(N^2 \log(N^2)) \implies \mathcal{O}(N^2 \log N)
\end{equation}

Per quanto riguarda la memoria, la ricorsione viene implementata con funzioni standard C++. L'occupazione di memoria totale per il riordinamento rimane quindi limitata a $\mathcal{O}(N^2)$.

\vspace{0.5cm}

\subsection{Modulo 3: Assemblaggio del Sistema in $\mathcal{O}(N^2)$}
La discretizzazione differenziale si basa su uno stencil discreto di Laplace a 5 punti. Il cuore matematico è attraversato serialmente sulle $N^2$ righe ideali che comporranno la matrice finale di Poisson.

Per la struttura fisica strettamente spaziale intrinseca all'operatore di Laplace bidimensionale $\nabla^2$, ciascuna molecola discreta si accoppia e interagisce analiticamente solo coi suoi diretti vicini (Nord, Sud, Est, Ovest). Il calcolatore impiega calcoli in esecuzione rapida $\mathcal{O}(1)$ per derivare i divisori spaziali $\kappa / h^2$ e incorporare con il segno opposto i contributi asimmetrici delle condizioni di Dirichlet al bordo $u_0(x,y)$.

Questa metodologia partorisce una matrice \textit{sparsa} che assicura a livello algebrico la genesi di non più di $5$ elementi non nulli per ciascuna riga del dominio. I moduli precludono alla radice l'inizializzazione di pesanti matrici dense temporanee (evitando allocazioni inutili di zeri); ragion per cui il tempo di saturazione dell'assemblaggio evolve in proporzione esatta a $V$, ovverosia $\mathcal{O}(N^2)$.
Infine, per abbattere criticamente i costi accessori dettati dalle riallocazioni dinamiche in volo, il vettore di triplette (formato COO) proprio della \textit{Standard Library} del C++ pre-riserva la dimensione necessaria impiegando la macro-chiamata \texttt{reserve(5*N*N)}. Questa istruzione chirurgica inchioda l'impronta di RAM a una mole rigidamente stimata in $\mathcal{O}(N^2)$.
% ------------------------------------------------------------
\newpage
\section{Piano di validazione sperimentale}
% ------------------------------------------------------------

L'obiettivo di questa fase di validazione architetturale è garantire la totale solidità dello scambio dati tra il backend C++ e l'interfaccia Python, nonché assicurare che la formulazione numerica del sistema lineare rispetti rigorosamente i principi fisici del problema. L'analisi approfondita delle performance computazionali (riduzione del fill-in e tempi di esecuzione) è invece demandata e ampiamente discussa all'interno del Jupyter Notebook.

\subsection*{1. Sanity Checks Fisici (Verifica della Soluzione)}

Il modo più efficace per validare l'intero stack (generazione della griglia, assemblaggio e risoluzione) consiste nel testare il sistema con scenari fisici elementari di cui si conosce a priori la soluzione analitica. Al fine di automatizzare questa batteria di test tramite la pipeline Python, \textbf{le condizioni al bordo $u_0(x,y)$ vengono lette dinamicamente tramite \texttt{stdin}} nel modulo \texttt{task3\_assemble.cpp}. Questo espediente permette di validare in modo istantaneo i seguenti scenari:

\begin{itemize}
  \item \textbf{Il "Caso Zero" (Opzione 0):} Impostando la sorgente interna $f = 0$ e la temperatura al bordo $u_0 = 0$, la soluzione numerica calcolata deve risultare identica a zero in ogni punto del dominio. L'aderenza a questo test esclude inequivocabilmente la presenza di perdite logiche nell'assemblaggio del termine noto $b$.
  \item \textbf{Caso Temperatura Costante (Opzione 1):} Fissando uniformemente l'involucro a una costante termica $C$ (es. $u_0 = 10$), in perfetta assenza di gradienti termici interni, la soluzione calcolata deve convergere univocamente a $C$ su tutto il dominio.
  \item \textbf{Caso Funzione Armonica (Opzione 4):} Attivando la condizione armonica quadratica $u_0 = x^2 - y^2$ con $f=0$, la soluzione numerica è forzata a ricalcare la medesima forma originaria in virtù del teorema della divergenza (\(\nabla^2(x^2 - y^2) = 0\)). Testando questa configurazione, lo scarto residuo della soluzione numerica si mantiene confinato nell'ordine della precisione di macchina ($\approx 10^{-16}$), confermando empiricamente la perfetta validità e accuratezza dello stencil a 5 punti.
\end{itemize}

\subsection*{2. Validazione e Integrità dello Scambio Dati}

Poiché l'architettura sfrutta esplicitamente il file system come veicolo di comunicazione tra i layer, è essenziale prevenire qualsiasi forma di corruzione dei dati in transito:
\begin{itemize}
  \item \textbf{Precisione Matematica Flottante in transito:} Per sventare l'introduzione di errori di arrotondamento durante il salvataggio dei coefficienti del sistema lineare, i flussi di scrittura \texttt{std::ofstream} nel C++ vengono rigorosamente forzati a una risoluzione nominale elevata tramite \texttt{std::setprecision(10)}. In questo modo, il risolutore in Python eredita una rappresentazione fedele e non degradata del Laplaciano discreto e del termine noto, mantenendo intatte le proprietà della matrice (in primis la simmetria e la definitezza positiva).
\end{itemize}
% ============================================================
\end{document}
% ============================================================
